\documentclass{article}%
\usepackage[T1]{fontenc}%
\usepackage[utf8]{inputenc}%
\usepackage{lmodern}%
\usepackage{textcomp}%
\usepackage{lastpage}%
\usepackage{relsize}%
\usepackage{ragged2e}%
%

\usepackage[utf8]{inputenc}
\usepackage{amsmath}
\usepackage{hyperref}
\usepackage{array}
\usepackage{graphicx}
\usepackage{longtable}
\usepackage{tikz}
\usepackage{titlesec}
\usepackage{subcaption}
\usepackage[margin=1in]{geometry}
\usepackage{float}
\usepackage{multicol}
\usepackage{hyperref}
\usepackage{url}

\newcommand{\position}{z}
\newcommand{\research}{r}
\newcommand{\price}{p}
\newcommand{\observed}{o}
\newcommand{\regressionprice}{y}
\newcommand{\currenttime}{t}
\newcommand{\constantc}{c}
\newcommand{\minimumratio}{K}
\newcommand{\lookbacklength}{L}
\newcommand{\dummytime}{\tau}
\newcommand{\dummyiterator}{n}
\newcommand{\dummyiteratortwo}{m}
\newcommand{\dummyiteratorthree}{l}
\newcommand{\decaycoefficient}{\phi}
\newcommand{\decaycoefficientone}{\decaycoefficient_{1}}
\newcommand{\decaycoefficienttwo}{\decaycoefficient_{2}}
\newcommand{\decaycoefficientthree}{\decaycoefficient_{3}}
\newcommand{\numberofcoefficients}{\Phi}
\newcommand{\fixedconstant}{\Omega}
\newcommand{\amplitudecoefficient}{\kappa}
\newcommand{\amplitudecoefficientone}{\amplitudecoefficient_{1}}
\newcommand{\amplitudecoefficienttwo}{\amplitudecoefficient_{2}}
\newcommand{\amplitudecoefficientthree}{\amplitudecoefficient_{3}}
\newcommand{\horizon}{h}
\newcommand{\contribution}{\lambda}
\newcommand{\contributionone}{\contribution_{1}}
\newcommand{\contributiontwo}{\contribution_{2}}
\newcommand{\bigcontribution}{\Lambda}
\newcommand{\bigcontributionone}{\bigcontribution_{1}}
\newcommand{\bigcontributiontwo}{\bigcontribution_{2}}
\newcommand{\bigcontributionthree}{\bigcontribution_{3}}
\newcommand{\bigcontributionshort}{\bigcontribution(\currenttime, \averagelengthshort, \price)}
\newcommand{\bigcontributionlong}{\bigcontribution(\currenttime, \averagelengthlong, \price)}
\newcommand{\bigcontributiontime}{\bigcontribution_\currenttime}
\newcommand{\averagelength}{L}
\newcommand{\averagelengthshort}{\averagelength_{1}}
\newcommand{\averagelengthlong}{\averagelength_{2}}
\newcommand{\kellyfraction}{F}
\newcommand{\rmserrorsquared}{\varepsilon^{{2}}_{rms_{\currenttime}}}
\newcommand{\rmserror}{\varepsilon_{rms_{\currenttime}}}
\newcommand{\amplitudecoefficientshort}{\amplitudecoefficientone}
\newcommand{\amplitudecoefficientlong}{\amplitudecoefficienttwo}
\newcommand{\genericfunction}{\zeta}
\newcommand{\scaleratiofactor}{\psi}
\newcommand{\Figure}[1]{Figure #1}
\newcommand{\tblofcontents}{\tableofcontents \pagebreak}
\newcommand{\portfolio}{V}

\titleformat*{\section}{\huge\bfseries}

\newcommand{\logo}
[0]
{\begin{figure}\centering\scalebox{0.25}{\includegraphics{../images/logo.png}}\end{figure}}
\newcommand{\rulename}[1]
{\section*{Predictive Relationship: #1}}
\newcommand{\ruledescription}[1]{\hspace{200mm}\section{Trading Strategy Description}#1}
\newcommand{\relationshipdescription}[1]{\hspace{200mm}\section{Relationship Description}#1}

%Tried to make recursive function. Table lines are buggy
\makeatletter
\newcommand{\ruleparameters}{\section{Rule Parameters}
Below is a table summarizing the parameters specific to this trading rule.
\begin{center}
\begin{tabular}{|m{9em}| m{3em}| m{15em}| m{3em}|}
\hline
Parameter Name & Default Value &  Description & Symbol
\@tablefilli
}

\newcommand{\fixedruleparameters}{\section{Fixed Strategy Parameters}
Below is a table summarizing the parameters specific to this trading rule.
\begin{center}
\begin{tabular}{|m{9em}| m{3em}| m{15em}| m{3em}|}
\hline
Parameter Name & Default Value &  Description & Symbol
\@tablefilli
}
\newcommand\@tablefilli{\@ifnextchar\stoptable{\@tablesend}{\@tablefillii}}
\newcommand\@tablefillii[4]{%
  \@tablefilliii{#1}{#2}{#3}{#4}
  \@tablefilli % restart the recursion
}
\newcommand\@tablefilliii[4]{%
  \\ \hline
  #1 & #2 & #3 & #4
}
\newcommand\@tablesend[1]{% The argument is \stoptable
  \\ \hline
  \end{tabular}
  \end{center}}
\makeatother


\newcommand{\furtherlinks}[0]{
\hspace{200mm}
\hspace{200mm}
\section*{Further Links}
\begin{enumerate}
    \item InferTrade: \url{https://www.infertrade.com}
    \item Privacy Policy/Legal notice: \url{https://www.infertrade.com/privacy-policy}
    \item InferStat Ltd: \url{https://www.inferstat.com}
\end{enumerate}
}


\newcommand{\appendixinfertrade}[0]{
    \hspace{200mm}
    \hspace{200mm}
    \section*{Appendix}
}


\newcommand{\regressionanderrorpartone}{
    \subsection*{Calculating Regression Error}
    \vspace{1mm}
    \justify This PDF omits any explanation of how the regression coefficients $k_1$ and $k_2$ are calculated as we are using a standard regression for calculation of both the coefficients and standard error. For InferTrade we used the open source \textit{sklearn} Python library. The same can be achieved in Microsoft Excel by using the \textit{FORECAST.LINEAR} function.

    \vspace{1mm}
    \justify As a worked illustration of the error calculation, assume that InferTrade predicted that a security would increase in price by $0.1\%$ between yesterday and today, and the price actually decreased by $4.9\%$. In this example, $\Delta \price_{\currenttime} = +0.1\%$, and $\Delta\observed_{\currenttime} = -4.9\%$, giving $\Delta P_{\currenttime} - \Delta\price_{\currenttime} = 0.1\% - 4.9\% = -5\% = -0.05$. This value is then squared to give $(\Delta\price_{\currenttime} - \Delta\observed_{\currenttime})^{2} = (-0.05)^2 = 0.0025$. To obtain $\rmserrorsquared$ for day $\currenttime$, this calculation is repeated for every day from day $i = \currenttime-L$ to day $i = \currenttime$, with values summed before dividing by the regression length $\lookbacklength$.

}

\newcommand{\regressionanderrorparttwo}{
    \vspace{1mm}


\begin{equation}
\Delta \price_{\currenttime+1} = k_1 \research_\currenttime+k_2 + \epsilon_{\currenttime+1}
\tag{\ref{eqn:levelregression} }
\end{equation}

 As $\lookbacklength = 120$ for all simplified mode relationships, this is the same as squaring the root-mean-square of the previous 120 prediction errors .
    \vspace{1mm}

    \justify This section can be skipped or skimmed by those comfortable with implementing a standard rolling linear regression and a standard rolling error. As seen previously, InferTrade uses a standard linear regression over a historical rolling window of length $\lookbacklength$ to construct an estimator for daily price changes in the form:
    

    
    \vspace{1mm}
    \justify \textbf{NB} For the first $L$ days, $\Delta\price_{\currenttime+1}$ calculates the in-sample prediction, meaning that the one regression model is calculated over the first $L$ data points and used for predictions on each of these days. For any subsequent days, $\Delta\price_{\currenttime+1}$ calculates the out-of-sample prediction, meaning that a new regression model is fit for each trailing window. 
    
    \vspace{1mm}
    \justify While $\Delta\price_{\currenttime+1}$ is the regression estimator used to predict future price changes with an observed error $\epsilon_\currenttime$, the regression itself is calculated using the actual observed price changes $\Delta\observed_\currenttime$. The standard root-mean-squared error factor, which scales the suggested allocation on day $\currenttime$, is squared to obtain $\rmserrorsquared$ and is calculated in percentage terms as follows:
    
    \\ % creates some space after equation
    \vspace{1mm}
    \begin{equation}
        \rmserrorsquared = 
        \begin{cases}
        \frac{1}{\lookbacklength} \mathlarger{\sum}_{i=\currenttime-L}^{\currenttime}(\Delta \price_{i} - \Delta \observed_{i})^{2} & \text{if } \currenttime > L,\\
        0 & \text{otherwise}
        \end{cases}
    \end{equation}
    \vspace{1mm}
    \justify
    
    
    \justify where $\Delta\price_{i}$ and $\Delta\observed_{i}$ are the predicted and observed percentage price changes from day $i-1$ to day $i$, and $\lookbacklength$ is the regression length. 
    
    \\ % creates some space after equation
    \vspace{1mm}
    \justify For example, assume that InferTrade predicted that a security would increase in price by $0.1\%$ between yesterday and today, and the price actually decreased by $4.9\%$. In this example, $\Delta \price_{\currenttime} = +0.1\%$, and $\Delta\observed_{\currenttime} = -4.9\%$, giving $\Delta P_{\currenttime} - \Delta\price_{\currenttime} = 0.1\% - 4.9\% = -5\% = -0.05$. This value is then squared to give $(\Delta\price_{\currenttime} - \Delta\observed_{\currenttime})^{2} = (-0.05)^2 = 0.0025$. To obtain $\rmserrorsquared$ for day $\currenttime$, this calculation is repeated for every day from day $i = \currenttime-L$ to day $i = \currenttime$, with values summed before dividing by the regression length $\lookbacklength$. As $\lookbacklength = 120$ for all simplified mode relationships, this is the same as squaring the root-mean-square of the previous 120 prediction errors .
    \vspace{1mm}
    \justify Note that if the magnitude $|\rmserror| < 1$, dividing by $\rmserrorsquared$ increases the suggested allocation. This has the effect of suggesting higher allocations when the rolling error is small, thereby providing some adjustment for risk arising from volatility in the regression estimate. If, in the example above, InferTrade predicts that tomorrow's price will increase by $0.1\%$, and the root-mean-square of the previous 120 errors $(\Delta\price_{i} - \Delta\observed_{i})$ is found to be $5\%$, then $\rmserrorsquared = 0.05^2 = 0.0025$, and the suggested allocation $\position_\currenttime$ will be scaled by factor of $\frac{1}{0.0025} = 400$, giving
    
    \\ % creates some space after equation
    \vspace{1mm}
    \begin{equation}
    \position_\currenttime = \frac{E[\Delta \price_{\currenttime+1}]}{\rmserrorsquared} = \frac{0.001}{0.0025} = 0.4 = 40 \%
    \end{equation}
    
    \vspace{1mm}
    \justify $E[\Delta\price_{\currenttime+1}]$ should be read as "The Expected Value of $\Delta\price_{\currenttime+1}$" and signifies that the actual price change $\Delta o_{\currenttime+1}$ is likely to deviate from $\Delta\price_{\currenttime+1}$ with prediction error $\epsilon_{\currenttime+1}$.
}


\newcommand{\howtotrade}[2]{
\section{How to Trade}
In order to trade with the rules InferTrade provides, we calculate allocations for each day. We then allocate that fraction of our total portfolio value (cash and securities) to the market we are trading - to do this we buy or sell securities to reach the target allocation.
#2
\subsection*{How Allocation Determines Trade Size}
The allocation is the fractional amount of the portfolios value used to determine the size of the trading position. For example, if the allocation for Microsoft (MSFT) shares is 50\%, and we have \$100, we invest \$50 so that the value of held stock is the same as the value of held cash.
\subsection*{Rule Specific Trading Details}
#1}
\newcommand{\graphdir}[1]{../rule_graphs/#1}
\newcommand{\ruleperformance}[1]{
\section{Rule Performance Considerations}
#1
}
\newcommand{\assumptions}[0]{}
\newcommand{\keyterms}{
\section{Glossary}
\begin{itemize}
    \item \textbf{Bullish:} Positive outlook on the market. Expectation of positive returns.
    \item \textbf{Bearish:} Negative outlook on the market. Expectation of negative returns.
    \item \textbf{Allocation:} The allocation is the fractional amount of the portfolios value used to determine the size of the trading position.
    \item \textbf{Parameter:} Value used by the trading rule in the calculation for trading position
    \item \textbf{Trading Rule:} Strategy to determine when to buy, hold or sell a position.
\end{itemize}
}

%
\usepackage{tocloft}%
\renewcommand{\cfttoctitlefont}{\Large\bfseries}%
%
\begin{document}%
\normalsize%
\logo%
\rulename{Change Relationship}%
\tableofcontents%

\relationshipdescription{
\justify
\Centering\textbf{In a change relationship, recent changes in the current value of the signal directly affects future changes in the price:}

\begin{equation}
\mathlarger{
    \Delta \price_{\currenttimet+1} \propto k \Delta S_t + c
    }
\end{equation}
\justify where ~$\Delta \price_{t+1}$ represents the future percentage price change, $\Delta S_t$ represents the previous percentage change in the research signal at current time $t$, $k$ is a scalar, and $c$ is a constant. 
\justify Intuitively, assuming a positive coefficient for k,  this means that when the signal is high it is good to buy the market, as the price will tend to drift higher. When the signal is low, it is good to sell the market as the price will tend to drift lower.\par
\vspace{1mm}
\justify
\Centering An example of a time series that might show a predictive change relationship would be interest rate:
\justify\Centering\textit{interest rate falls -> recent change in interest rate is high -> price of the S&P 500 goes up because investors are allocating capital towards higher return generating assets}
\vspace{1mm}
\justify
We could also see a change component to price forecasts or fair value models if an analyst or model is updated infrequently, as a recent positive change could indicate incorporation of fresh news or information. However if the market price has moved stably then the new information may already be fully incorporated in the price.
\vspace{1mm}
\justify
Similarly, a technical positioning indicator as a signal may be more likely to show a change relationship, whereby changes in the value of a signal drive the price, because a stable level of interest may not move the market. We would typically not expect a signal that is a price forecast to show a predictive change relationship, as a price forecast is typically relative to the current price (a difference relationship).
}%
\ruledescription{
\justify

A predictive change relationship can be reflected in many kinds of rules. InferTrade uses a 120 period (6 months for daily data) rolling regression of the percentage change in the signal from the prior day against next day's price change as a benchmark. This trading rule recommends taking a position sized as a percentage of the capital you have allocated to trading this particular market.
\vspace{1mm}
\justify
This rule will show higher returns than usual after optimisation if a significant change relationship is present between the price and signal series. The following equation shows how a change regression trading rule calculates a position sizing:

\begin{equation}
\position_\currenttime = 
k_{1}\left(\frac{\research_t}{\research_{t-1}}-1\right)+k_2 + \epsilon_{\currenttime+1}
\end{equation}
\\ % creates some space after equation
where $\position_\currenttime$ is the recommended strategy allocation at time $\currenttime$, $\research_t$ is the research value at time $\currenttime$, $k_1$ is the change coefficient and $k_2$ is the static coefficient. The bracketed term calculates the percentage change in the value of the research signal from time $t$ to $t-1$.
} 

\justify with the error symmetric such that 

\begin{equation}
E^{model}_{\currenttime}[\Delta \price_\currenttime+1] =k_{1}\left(\frac{\research_t}{\research_{t-1}}-1\right)+k_2 
\label{eqn:regressionforecast}
\end{equation}%Argument is rule description
\justify where $E^{model}_{\currenttime}[\Delta \price_{\currenttime+1}]$ is the predicted next day percentage price change, $\research_\currenttime$ is the research value at time $\currenttime$, and $k_1$ and $k_2$ are the change  and static coefficients. $\epsilon_\currenttime$ represents the in sample regression error between the model predicted and observed price change for time \currenttime, such that $\epsilon_{\currenttime+1}= E^{model}_{\currenttime-1}[\Delta \price_{\currenttime}] - \Delta \price_{\currenttime}$.
\vspace{1mm}
\justify To give a suggested portfolio allocation $\position_\currenttime$ for the market, e.g. ``invest 5\% of strategy allocation into the security'', the regression forecast's expected next day \% price move is divided by the square of the Standard Root Mean Squared Percentage Prediction Error $\rmserrorsquared$ in our above regression:
\vspace{1mm}
\\



\vspace{1mm}
\begin{equation}
    \position_\currenttime = 
    \begin{cases}
    \frac{\Delta \price_{\currenttime+1}}{\rmserrorsquared} & \text{if } \currenttime > L,\\
    0                                        & \text{otherwise}
    \end{cases}
\end{equation}



\vspace{1mm}
\justify where $\position_\currenttime$ is the suggested allocation as a percentage of the capital which you have allocated to trading this market at time $\currenttime$.
\justify Worked example for liquid market:  current share price is \$50, regression coefficients are 0.00005 for change and 0.0 for the constant, with \$1,000 allocated to this single-stock strategy.  The current research series - the interest rate - is +5, with a RMS prediction error from the regression of 5\%. In this situation the strategy recipe recommends an investment of $0.00005 * 5 / (0.05)^2 = 10\%$ be allocated. So \$1,000 * 10\% = \$100 would be allocated, so we would buy $\frac{\$100}{\$100} = 2$ shares.
We see from this example the expected gain over the next day is 2.5 basis points (minus trading costs), versus 500 basis points of daily risk, so in practice this kind of strategy is most useful where the signal and price both has a predictive edge and is stable (auto-correlated) in recommendation over time, so that the `edge' from multiple time steps can compound linearly without incurring additional bid-offer spreads whilst the `risk' aggregates scales with the square-root of time (for random walk).


\vspace{1mm}
\justify (Note that if the magnitude of the error term is smaller than the magnitude of the price change prediction, i.e. $|\rmserror| < |\Delta\price_{\currenttime+1}|$, dividing $\Delta\price_{\currenttime+1}$ by $\rmserrorsquared$ would result in $\position_\currenttime>1$. Limits can then be applied to cap the allocation. E.g. between 0\% and 100\% if short selling is prohibited. In practice it is generally unlikely that in liquid markets a genuine statistical trading signal this strong will be found.)

\vspace{1mm}
\ruleparameters
{Regression Length}{120}{This is the number of previous days used to estimate the regression coefficients.}{$\lookbacklength$}%
\stoptable 

\vspace{1mm}
\justify Note that rule parameters are fixed to their default values for this and any other relationships optimised on simplified mode. For estimating the regression coefficients, a rolling window is used, so that at time t, the regression data uses research data from time t-121 to t-1 to calculate 120 research changes as the regression's independent variable. Price data from t-120 to t is used to calculate 120 next day price changes as the regression's dependent variable. %Argument is rule description
\assumptions%
\keyterms%
\furtherlinks%
\appendixinfertrade{
\vspace{1mm}
\regressionanderrorpartone}
\end{document}
